%% parts/p04/complexity.tex
%% Registered new content D-P5-007/N-a-064, with apparatus A-050/A-052/A-127.
%% Hierarchy and completeness results are citation-only.

\chapter{Computational Complexity}
\label{sec:p04-complexity}

\noindent\textit{Reader note.} This separately authored chapter extends the
exact inherited computability sequence with standard complexity and
descriptive-complexity results.  Its machine conventions, class definitions,
and cited theorems do not replace any inherited source unit.

Computability asks whether some machine eventually decides a language.
Complexity also fixes a machine model, an encoding, and a resource measured as
a function of input length.  Formal-model equivalence therefore does not imply
equality of time or space costs.  The classes below use a standard multitape
Turing-machine model and the exact conventions in
\cite[chaps.~1--5]{AroraBarak2009}.

\section{Resource bounds, reductions, and classes}
\label{sec:p04-complexity-classes}

\begin{hbconvention}{def:p04-complexity-convention}{Complexity and reduction convention}
A decision problem is a language \(L\subseteq\{0,1\}^*\), and
\(n=\operatorname{len}(x)\) is the bit length of an input.  Machines have a read-only
input and finitely many work tapes.  For functions \(f,g:\N\to\N\), write
\(f(n)=O(g(n))\) iff there are \(c,n_0\in\N\), with \(c\geq1\), such that
\(f(n)\leq c g(n)\) for every \(n\geq n_0\); write \(f(n)=o(g(n))\) iff,
for every \(c\in\N\) with \(c\geq1\), there is \(n_0\in\N\) such that
\(c f(n)\leq g(n)\) for every \(n\geq n_0\).  Every logarithm in this
chapter has base \(2\).  A nondeterministic machine replaces the transition
function by a finite transition relation.  It accepts an input iff some
computation branch accepts, rejects iff every branch rejects, and counts as a
decider only when every branch halts; its time and work-tape space are the
maxima over all branches.  Unless another reduction is explicitly
named, \(A\leq_m^pB\) means that a deterministic polynomial-time computable
map \(f\) satisfies
\[
 x\in A\quad\Leftrightarrow\quad f(x)\in B.
\]
For a class \(\mathcal C\), a language is \(\mathcal C\)-hard when every
language in \(\mathcal C\) reduces to it by \(\leq_m^p\), and is
\(\mathcal C\)-complete when it is also a member of \(\mathcal C\).
Membership and hardness are both required; ``complete'' never means
deductively or semantically complete in this chapter.
\index{polynomial-time many-one reduction}
\index{complexity class!completeness}
\index{asymptotic notation}
\end{hbconvention}

\begin{hbdefinition}{def:p04-resource-classes}{Time and space classes}
For a function \(f:\N\to\N\), \(\mathsf{DTIME}(f(n))\) and
\(\mathsf{NTIME}(f(n))\) are the languages decided within \(O(f(n))\) steps
by deterministic and nondeterministic machines, respectively.
Similarly, \(\mathsf{SPACE}(f(n))\) and \(\mathsf{NSPACE}(f(n))\) measure the
number of work-tape cells used; the read-only input is not counted.  Put
\[
\begin{aligned}
 \mathsf P&=\bigcup_{k\geq1}\mathsf{DTIME}(n^k),&
 \mathsf{NP}&=\bigcup_{k\geq1}\mathsf{NTIME}(n^k),\\
 \mathsf{coNP}&=\{\{0,1\}^*\setminus L:L\in\mathsf{NP}\},&
 \mathsf{PSPACE}&=\bigcup_{k\geq1}\mathsf{SPACE}(n^k),\\
 \mathsf{NPSPACE}&=\bigcup_{k\geq1}\mathsf{NSPACE}(n^k),&
 \mathsf{EXPTIME}&=\bigcup_{k\geq1}\mathsf{DTIME}(2^{n^k}),\\
 &&\mathsf{EXPSPACE}&=\bigcup_{k\geq1}\mathsf{SPACE}(2^{n^k}).
\end{aligned}
\]
These are classes of encoded languages.  The symbol \(\mathsf{coNP}\), for
example, names the complement class and not a claim that
\(\mathsf{NP}=\mathsf{coNP}\).
\index{P@\(\mathsf P\)}
\index{NP@\(\mathsf{NP}\)}
\index{PSPACE@\(\mathsf{PSPACE}\)}
\index{EXPTIME@\(\mathsf{EXPTIME}\)}
\index{EXPSPACE@\(\mathsf{EXPSPACE}\)}
\end{hbdefinition}

\begin{hbdefinition}{def:p04-constructible-bound}{Constructible resource bound}
A function \(f:\N\to\N\) with \(f(n)\geq n\) is \emph{time constructible} when a deterministic
machine, on input \(1^n\), can produce the unary clock value \(1^{f(n)}\) in
\(O(f(n))\) time.  A function with
\(f(n)\geq\lceil\log(n+1)\rceil\) is \emph{space constructible} when a deterministic
machine, on input \(1^n\), uses exactly \(f(n)\) work cells and halts.  These
hypotheses prevent a hierarchy statement from hiding a resource bound that a
machine cannot itself track.
\index{time-constructible function}
\index{space-constructible function}
\end{hbdefinition}

\begin{hbdefinition}{def:p04-polynomial-hierarchy}{Polynomial hierarchy}
Put \(\Sigma^p_0=\Pi^p_0=\mathsf P\).  For \(k\geq1\), a language \(L\)
belongs to \(\Sigma^p_k\) when there are a polynomial \(q\) and a
polynomial-time predicate \(R\) such that
\[
 x\in L\quad\Leftrightarrow\quad
 \begin{cases}
  \exists y_1\,\forall y_2\,\exists y_3\cdots\exists y_k\;
    R(x,y_1,\ldots,y_k),&k\text{ odd},\\
  \exists y_1\,\forall y_2\,\exists y_3\cdots\forall y_k\;
    R(x,y_1,\ldots,y_k),&k\text{ even},
 \end{cases}
 \qquad \operatorname{len}(y_i)\leq q(\operatorname{len}(x)).
\]
There are \(k\) alternating bounded string-quantifier blocks; the parity
clauses make the final quantifier explicit.  The class
\(\Pi^p_k\) is defined by the dual prefix beginning universally, and
\[
 \mathsf{PH}\defeq\bigcup_{k\geq0}\Sigma^p_k.
\]
The equivalent oracle definition is not needed as a second definition here;
its equivalence with this bounded-quantifier presentation is part of the
cited theory \cite{Stockmeyer1976,AroraBarak2009}.
\index{polynomial hierarchy}
\end{hbdefinition}

\section{Hierarchy and simulation theorems}
\label{sec:p04-complexity-hierarchies}

\begin{hbtheorem}{thm:p04-time-hierarchy}{cited=HartmanisStearns1965}{Deterministic time hierarchy theorem}
If \(f\) and \(g\) are time-constructible and
\(f(n)\log f(n)=o(g(n))\), then
\[
 \mathsf{DTIME}(f(n))\subsetneq\mathsf{DTIME}(g(n)).
\]
In particular, \(\mathsf P\subsetneq\mathsf{EXPTIME}\).  The
constructibility and simulation-overhead gap are hypotheses of the theorem,
not optional regularity assumptions \cite[thm.~3.1]{AroraBarak2009}.
\end{hbtheorem}

\begin{hbtheorem}{thm:p04-space-hierarchy}{cited=StearnsHartmanisLewis1965}{Deterministic space hierarchy theorem}
If \(f\) and \(g\) are space-constructible and \(f(n)=o(g(n))\), then
\[
 \mathsf{SPACE}(f(n))\subsetneq\mathsf{SPACE}(g(n)).
\]
In particular,
\(\mathsf{PSPACE}\subsetneq\mathsf{EXPSPACE}\)
\cite[thm.~3.2]{AroraBarak2009}.
\end{hbtheorem}

\begin{hbtheorem}{thm:p04-savitch}{cited=Savitch1970}{Savitch's theorem}
If \(f(n)\geq\log n\) is space constructible, then
\[
 \mathsf{NSPACE}(f(n))
 \subseteq \mathsf{SPACE}(f(n)^2).
\]
Consequently \(\mathsf{NPSPACE}=\mathsf{PSPACE}\).  The lower bound and
constructibility hypothesis remain part of the general theorem even though
the polynomial-space corollary suppresses them
\cite[thm.~4.12]{AroraBarak2009}.
\end{hbtheorem}

\begin{hbwarning}{def:p04-hierarchy-theorems-scope}{What the hierarchy theorems do not separate}
The preceding theorems give the global strict containments
\(\mathsf P\subsetneq\mathsf{EXPTIME}\) and
\(\mathsf{PSPACE}\subsetneq\mathsf{EXPSPACE}\).  They do not prove that
any adjacent inclusion in
\[
 \mathsf P\subseteq\mathsf{NP}\subseteq\mathsf{PH}
 \subseteq\mathsf{PSPACE}\subseteq\mathsf{EXPTIME}
 \subseteq\mathsf{EXPSPACE}
\]
is strict.
\end{hbwarning}

\section{Logic-relevant complete problems}
\label{sec:p04-complexity-complete-problems}

The SAT and TAUT languages and their fixed binary encoding were defined in
Definition~\ref{def:p02-sat-taut}; malformed strings remain no instances.

\begin{hbdefinition}{def:p04-qbf-languages}{Quantified Boolean languages}
A \emph{quantified Boolean formula} (QBF) has an explicitly encoded propositional
matrix and quantifiers ranging over \(\{0,1\}\).  It is \emph{closed} when it
has no free propositional variable.  The language of true quantified Boolean
formulas, \(\mathsf{TQBF}\), contains
exactly the encodings of true closed quantified Boolean formulas, with no
bound on alternation.

For fixed \(k\geq1\), \(\Sigma_k\text{-}\mathsf{SAT}\) contains true closed
prenex formulas with \(k\) alternating quantifier blocks beginning
existentially.  The dual language \(\Pi_k\text{-}\mathsf{SAT}\) begins
universally.  These fixed-alternation languages are distinct from
\(\mathsf{TQBF}\), which permits unbounded alternation.
\index{quantified Boolean formula}
\index{TQBF}
\end{hbdefinition}

\begin{hbdefinition}{def:p04-succinct-problems}{Circuit value and succinct encodings}
The \emph{circuit value problem} \(\mathsf{CVP}\) takes an acyclic Boolean
circuit with constants and \(\neg,\land,\lor\) gates in topological order and
asks whether its designated output gate evaluates to \(1\).  Write
\(A\leq_m^{\log}B\) when the many-one translator is computed by a deterministic
logarithmic-work-space transducer.

A succinct instance uses a fixed self-delimiting gate-list or QBF encoding
whose decoder ignores trailing zero padding after its end marker.  Validity and
decoding are logspace computable, and the padded convention is logspace
intertranslatable with the ordinary encodings used above.  Thus every ordinary
instance has valid encodings of arbitrarily large power-of-two length.

A succinct instance is a Boolean circuit \(S\) with \(m\) input bits.  It
describes the \(2^m\)-bit string \(z_S\) whose bit at position \(i\) is
\(S(\operatorname{bin}_m(i))\), where \(\operatorname{bin}_m(i)\) is the
unique length-\(m\) binary expansion of \(0\leq i<2^m\), padded with leading
zeros.  This fixed-width address differs from the variable-length canonical
numeral encoding used for machine inputs.  The \emph{succinct circuit value problem}
\(\mathsf{SCVP}\) asks whether \(z_S\) is a valid encoded circuit with output
value \(1\).  The \emph{succinct TQBF problem} \(\mathsf{STQBF}\) asks whether
\(z_S\) is a valid encoded true closed quantified Boolean formula.  Invalid
decoded strings are no instances.

Succinctness is needed to present an exponentially large ordinary instance by
a polynomial-size circuit.  It is an encoding convention, not a new machine
resource or logical connective.  It differs from ordinary CVP/TQBF by the
exponential decompression and is used to calibrate the two exponential classes
in the completeness catalog.
\index{circuit value problem}\index{succinct encoding}
\end{hbdefinition}

\begin{hbtheorem}{thm:p04-succinct-completeness}{cited=Papadimitriou1994}{Canonical complete problems for \(\mathsf P\), \(\mathsf{EXPTIME}\), and \(\mathsf{EXPSPACE}\)}
The language \(\mathsf{CVP}\) is \(\mathsf P\)-complete under
logspace many-one reductions.  The language \(\mathsf{SCVP}\) is
\(\mathsf{EXPTIME}\)-complete and \(\mathsf{STQBF}\) is
\(\mathsf{EXPSPACE}\)-complete under polynomial-time many-one reductions.
\end{hbtheorem}

\begin{hbtheorem}{thm:p04-cook-levin}{cited=AroraBarak2009}{Cook--Levin theorem}
The language \(\mathsf{SAT}\) belongs to \(\mathsf{NP}\), and every language
in \(\mathsf{NP}\) has a polynomial-time many-one reduction to
\(\mathsf{SAT}\).  Hence \(\mathsf{SAT}\) is \(\mathsf{NP}\)-complete under
the reduction in Convention~\ref{def:p04-complexity-convention}; this exact
modern encoding-level formulation is \cite[thm.~2.10]{AroraBarak2009}.
Historically, \textcite{Cook1971} proves the corresponding tautology result
using polynomial-time query-machine \(P\)-reducibility, rather than the
many-one convention fixed here; \textcite{Levin1973} is the independent
historical source.
\end{hbtheorem}

\begin{hbtheorem}{thm:p04-tqbf-pspace}{proved=annex:p9.tqbf.1}{TQBF completeness}
The language \(\mathsf{TQBF}\) is \(\mathsf{PSPACE}\)-complete under
polynomial-time many-one reductions.  Both membership and hardness are part
of the statement.  The historical source role is
\textcite{StockmeyerMeyer1973}; a complete handbook proof is
Annex~\ref{annex:p9.tqbf.1} \cite[thm.~4.11]{AroraBarak2009}.
\end{hbtheorem}

\begin{hbtheorem}{thm:p04-fixed-qbf-completeness}{cited=Stockmeyer1976}{Fixed-alternation QBF completeness}
For each fixed \(k\geq1\),
\(\Sigma_k\text{-}\mathsf{SAT}\) is \(\Sigma^p_k\)-complete and
\(\Pi_k\text{-}\mathsf{SAT}\) is \(\Pi^p_k\)-complete under
polynomial-time many-one reductions.
\end{hbtheorem}

\begingroup\small
\begin{longtable}{@{}p{0.17\textwidth}p{0.42\textwidth}p{0.32\textwidth}@{}}
\caption{Logic-relevant completeness catalog, paired with the class ladder in
Table~\ref{tab:p04-complexity-ladder}.}
\label{tab:p04-completeness-catalog}\\
\toprule
Problem & Encoded yes instances & Exact classification and evidence \\
\midrule
\endfirsthead
\toprule
Problem & Encoded yes instances & Exact classification and evidence \\
\midrule
\endhead
\bottomrule
\endfoot
\(\mathsf{CVP}\) & Valid encoded acyclic Boolean circuits whose designated
output evaluates to \(1\). & \(\mathsf P\)-complete under logspace many-one
reductions; Theorem~\ref{thm:p04-succinct-completeness}. \\
\addlinespace
\(\mathsf{SAT}\) & Satisfiable encoded propositional formulas under
Definition~\ref{def:p02-sat-taut}. & \(\mathsf{NP}\)-complete;
Theorem~\ref{thm:p04-cook-levin}; \cite[thm.~2.10]{AroraBarak2009}. \\
\addlinespace
\(3\mathsf{SAT}\) & Satisfiable encoded 3-CNF formulas, each clause having at
most three literals. & \(\mathsf{NP}\)-complete under the same reductions;
\cite[thm.~2.10]{AroraBarak2009}; \cite{Karp1972}. \\
\addlinespace
\(\mathsf{TAUT}\) & Encoded propositional formulas true under every Boolean
assignment. & \(\mathsf{coNP}\)-complete;
Theorem~\ref{thm:p02-taut-conp-complete}; \cite{AroraBarak2009}. \\
\addlinespace
\(\Sigma_k\text{-}\mathsf{SAT}\) & True closed prenex QBFs with fixed
\(k\geq1\), beginning existentially. & \(\Sigma^p_k\)-complete;
Theorem~\ref{thm:p04-fixed-qbf-completeness}; \cite{Stockmeyer1976}. \\
\addlinespace
\(\Pi_k\text{-}\mathsf{SAT}\) & The dual fixed-\(k\) prefix beginning
universally. & \(\Pi^p_k\)-complete;
Theorem~\ref{thm:p04-fixed-qbf-completeness}; \cite{Stockmeyer1976}. \\
\addlinespace
\(\mathsf{TQBF}\) & True closed QBFs with unbounded alternation. &
\(\mathsf{PSPACE}\)-complete; Theorem~\ref{thm:p04-tqbf-pspace};
\cite{StockmeyerMeyer1973}. \\
\addlinespace
\(\mathsf{SCVP}\) & Succinct descriptions of exponentially large valid
circuits evaluating to \(1\). & \(\mathsf{EXPTIME}\)-complete;
Theorem~\ref{thm:p04-succinct-completeness}. \\
\addlinespace
\(\mathsf{STQBF}\) & Succinct descriptions of exponentially long valid true
closed QBFs. & \(\mathsf{EXPSPACE}\)-complete;
Theorem~\ref{thm:p04-succinct-completeness}. \\
\addlinespace
First-order validity & Codes of first-order sentences valid in every
structure under the fixed signature convention. & c.e. and undecidable, not
a finite-level completeness claim; Theorem~\ref{thm:fol-validity-undecidable}
and Corollary~\ref{thm:fol-validity-enumerable};
\cite{Church1936,Turing1937}. \\
\end{longtable}
\endgroup

\section{The class ladder and its open edges}
\label{sec:p04-complexity-ladder}

\begingroup\small
\begin{longtable}{@{}p{0.25\textwidth}p{0.25\textwidth}p{0.41\textwidth}@{}}
\caption{Complexity-class ladder.  Every row distinguishes
inclusion, equality, known strictness, and an unresolved equality.}
\label{tab:p04-complexity-ladder}\\
\toprule
Relation & Classification & Evidence and unresolved boundary \\
\midrule
\endfirsthead
\toprule
Relation & Classification & Evidence and unresolved boundary \\
\midrule
\endhead
\bottomrule
\endfoot
\(\mathsf P\subseteq\mathsf{NP}\) & proved inclusion & Immediate from the
machine definitions; equality is open \cite{AroraBarak2009}. \\
\addlinespace
\(\mathsf P\subseteq\mathsf{coNP}\) & proved inclusion & Deterministic
polynomial time is closed under complement; equality is open
\cite{AroraBarak2009}. \\
\addlinespace
\(\mathsf{NP},\mathsf{coNP}\subseteq\mathsf{PH}\) & proved inclusions &
\(\Sigma^p_1=\mathsf{NP}\) and \(\Pi^p_1=\mathsf{coNP}\);
\(\mathsf{NP}=\mathsf{coNP}\) and collapse of \(\mathsf{PH}\) are open
\cite{Stockmeyer1976}. \\
\addlinespace
\(\mathsf{PH}\subseteq\mathsf{PSPACE}\) & proved inclusion & Polynomially
many bounded alternations can be evaluated in polynomial space; equality and
strictness are open \cite[ch.~5]{AroraBarak2009}. \\
\addlinespace
\(\mathsf{PSPACE}=\mathsf{NPSPACE}\) & proved equality & Savitch's theorem,
Theorem~\ref{thm:p04-savitch}. \\
\addlinespace
\(\mathsf{PSPACE}\subseteq\mathsf{EXPTIME}\) & proved inclusion & A
polynomial-space computation has at most exponentially many configurations;
equality and strictness are open \cite[ch.~4]{AroraBarak2009}. \\
\addlinespace
\(\mathsf{P}\subsetneq\mathsf{EXPTIME}\) & proved strict containment &
Deterministic time hierarchy, Theorem~\ref{thm:p04-time-hierarchy}. \\
\addlinespace
\(\mathsf{EXPTIME}\subseteq\mathsf{EXPSPACE}\) & proved inclusion & A
machine uses no more work cells than elapsed steps; adjacent strictness is not
asserted \cite{AroraBarak2009}. \\
\addlinespace
\(\mathsf{PSPACE}\subsetneq\mathsf{EXPSPACE}\) & proved strict containment &
Deterministic space hierarchy, Theorem~\ref{thm:p04-space-hierarchy}. \\
\end{longtable}
\endgroup

\begin{figure}[htbp]
\centering
\begin{tikzpicture}[
  node distance=12mm and 18mm,
  cls/.style={draw,rounded corners,minimum width=24mm,minimum height=7mm,
              align=center,font=\small},
  inc/.style={-{Latex[length=2mm]},thick},
  strict/.style={-{Latex[length=2.4mm]},very thick},
  open/.style={densely dashed}
]
 \node[cls] (p) {\(\mathsf P\)};
 \node[cls,above left=of p] (np) {\(\mathsf{NP}\)};
 \node[cls,above right=of p] (conp) {\(\mathsf{coNP}\)};
 \node[cls,above=26mm of p] (ph) {\(\mathsf{PH}\)};
 \node[cls,above=of ph] (pspace) {\(\mathsf{PSPACE}=\mathsf{NPSPACE}\)};
 \node[cls,above=of pspace] (exptime) {\(\mathsf{EXPTIME}\)};
 \node[cls,above=of exptime] (expspace) {\(\mathsf{EXPSPACE}\)};
 \draw[inc] (p) -- (np);
 \draw[inc] (p) -- (conp);
 \draw[inc] (np) -- (ph);
 \draw[inc] (conp) -- (ph);
 \draw[inc] (ph) -- (pspace);
 \draw[inc] (pspace) -- (exptime);
 \draw[inc] (exptime) -- (expspace);
 \draw[open] (np) -- node[above,font=\scriptsize] {\(?=\)} (conp);
 \draw[open,bend left=18] (p) to node[left,font=\scriptsize] {\(?=\)} (np);
 \draw[open,bend right=18] (p) to node[right,font=\scriptsize] {\(?=\)} (conp);
 \draw[open,bend left=14] (np) to node[left,font=\scriptsize] {\(?=\)} (ph);
 \draw[open,bend right=14] (conp) to node[right,font=\scriptsize] {\(?=\)} (ph);
 \draw[open,bend right=26] (ph) to node[right,font=\scriptsize] {\(?=\)} (pspace);
 \draw[open,bend right=25] (pspace) to node[right,font=\scriptsize] {\(?=\)} (exptime);
 \draw[open,bend right=24] (exptime) to node[right,font=\scriptsize] {\(?=\)} (expspace);
 \draw[strict,bend left=44] (p) to node[left,font=\scriptsize,align=center]
   {known\\proper} (exptime);
 \draw[strict,bend left=42] (pspace) to node[left,font=\scriptsize,align=center]
   {known\\proper} (expspace);
\end{tikzpicture}
\caption{Complexity inclusions and open equality questions.
Solid arrows are proved inclusions; heavy curved arrows are the two proved
global strict containments; dashed connectors marked \(?=\) record unresolved
equalities.  The drawing is Hasse-style layout only: no edge is asserted to be
a cover relation.  Evidence is in
Table~\ref{tab:p04-complexity-ladder} and
Theorems~\ref{thm:p04-time-hierarchy}--\ref{thm:p04-savitch}.}
\label{fig:p04-complexity-inclusions}
\end{figure}

\section{Descriptive-complexity correspondences}
\label{sec:p04-descriptive-complexity}

\begin{hbdefinition}{def:p04-captures-complexity}{Capturing a class on finite structures}
Fix a finite relational vocabulary \(R_1,\ldots,R_r\), including its displayed
order and arities.  For \(n\geq1\), present an \(n\)-element structure on
\([n]=\{0,\ldots,n-1\}\).  If the standard binary expansion of \(n\) is the
\(k\)-bit word \(b\), its encoding begins with \(1^k0b\), followed, for
\(i=1,\ldots,r\), by the characteristic bits of \(R_i\) on
\([n]^{\operatorname{arity}(R_i)}\) in lexicographic tuple order.  An
arbitrarily presented finite structure is represented by each bijective
relabeling onto \([n]\); an isomorphism-invariant class therefore gives the
same membership answer to every such representation.

A logic formalism \(\mathfrak L\) \emph{captures} a complexity class
\(\mathcal C\) on the stated structures when:
\begin{enumerate}
 \item the encoded model class of every \(\mathfrak L\)-sentence belongs to
       \(\mathcal C\); and
 \item every isomorphism-invariant property whose encoded membership language
       belongs to \(\mathcal C\) is defined by a \(\mathfrak L\)-sentence.
\end{enumerate}
This finite-structure correspondence is neither ordinary definability over
arbitrary structures nor a claim that formulas and machines are the same
objects.  A built-in total order is mentioned exactly when the correspondence
requires it.
\index{descriptive complexity}
\index{captures a complexity class}
\end{hbdefinition}

\begin{hbdefinition}{def:p04-fol-lfp}{First-order least fixed-point logic}
Let \(R\) be a fresh \(r\)-ary relation variable and let
\(\varphi(R,\bar x,\bar y)\) be first-order with every occurrence of \(R\)
positive, meaning under an even number of negations in negation normal form.
For parameters \(\bar b\) in a structure \(M\), it defines the monotone
operator
\[
 F_{\bar b}(S)
 =\{\bar a\in M^r:M\entails\varphi(S,\bar a,\bar b)\}
 \quad(S\subseteq M^r).
\]
The formula
\([\operatorname{lfp}_{R,\bar x}\varphi](\bar t)\) holds exactly when the
value of \(\bar t\) belongs to the least fixed point of \(F_{\bar b}\).  On a
finite structure that fixed point is reached by the iteration
\(\varnothing,F(\varnothing),F^2(\varnothing),\ldots\).

Positivity is needed to make the operator monotone and the least fixed point
well defined.  This is an object-language formula constructor with the
displayed finite-structure semantics, not an arbitrary recursion scheme or a
second-order quantifier.  It differs from first-order logic by allowing this
iterated relation definition and is used in the Immerman--Vardi theorem and
the correspondence table below.
\index{least fixed-point logic}\index{FO(LFP)}
\end{hbdefinition}

Fagin's theorem is stated once, at
Theorem~\ref{thm:p03-fagin}: over finite relational structures,
existential second-order logic captures \(\mathsf{NP}\), with no built-in
order required.  The present chapter uses that same public theorem as a
descriptive-complexity correspondence rather than creating a second identity.

\begin{hbtheorem}{thm:p04-immerman-vardi}{cited=Immerman1986}{Ordered fixed-point capture}
Over finite relational structures equipped with a built-in total order,
first-order logic with a least fixed-point operator captures \(\mathsf P\)
\cite{Vardi1982}.  The order hypothesis is part of the statement.
\end{hbtheorem}

In the following table, FO abbreviates first-order logic, ESO existential
second-order logic, and FO(LFP) first-order logic with a least fixed-point
operator.  The later warning uses BIT for the built-in binary bit predicate.

\begin{table}[htbp]
\centering
\small
\begin{tabular}{@{}p{.14\textwidth}p{.32\textwidth}p{.20\textwidth}p{.25\textwidth}@{}}
\toprule
Logic & Finite-structure convention & Captured class & Evidence \\
\midrule
FO & finite relational structures; baseline only & no unqualified
machine-class equality asserted & \cite{Immerman1999} \\
ESO & finite relational structures; no built-in order required &
\(\mathsf{NP}\) & Theorem~\ref{thm:p03-fagin}; \cite{Fagin1974} \\
FO(LFP) & finite relational structures with a built-in total order &
\(\mathsf P\) & Theorem~\ref{thm:p04-immerman-vardi};
\cite{Immerman1986,Vardi1982} \\
\bottomrule
\end{tabular}
\caption{Descriptive-complexity correspondences.  The
structure and order columns are hypotheses, not expository decoration.}
\label{tab:p04-descriptive-complexity}
\end{table}

\begin{hbwarning}{def:p04-fo-aczero}{No unqualified FO--circuit-class equation}
No equation between plain first-order logic and a circuit class is asserted
here.  Such a statement would have to fix ordered finite structures,
built-in numerical predicates such as order or BIT, and a uniformity notion.
The ESO/\(\mathsf{NP}\) and ordered FO(LFP)/\(\mathsf P\) correspondences in
Table~\ref{tab:p04-descriptive-complexity} do not need that additional claim.
\end{hbwarning}

\section{Further reading: historical sources and modern treatments}
\label{sec:p04-further-reading}
\markright{25.6.\ FURTHER READING AND SOURCES}

\subsection*{Historical and primary sources}

For formal-model equivalence and the first undecidability results, begin with
Church and Turing \cite{Church1936,Turing1937,Turing1937Lambda}.  The degree
and c.e.-set landmarks are Post, Kleene--Post, Friedberg, Muchnik, and Myhill
\cite{Post1944,KleenePost1954,Friedberg1957,Muchnik1956,Myhill1955}.
The analytical, information, and Diophantine frontiers begin with
\cite{Kleene1955Hierarchy,Kolmogorov1965,Chaitin1974,Chaitin1975,
DavisPutnamRobinson1961,Matiyasevich1970}.  For word problems, Busy Beaver,
and tilings, see \cite{Post1947,Markov1947,Novikov1955,Boone1959,
Rado1962BusyBeaver,Berger1966}.  The complexity and descriptive-complexity
landmarks include \cite{HartmanisStearns1965,StearnsHartmanisLewis1965,
Cook1971,Levin1973,StockmeyerMeyer1973,Fagin1974,Immerman1986}.

\subsection*{Modern systematic treatments}

For computability, degrees, and the analytical boundary, see
\cite{Rogers1967,Soare2016,Moschovakis2016Hyperarithmetical}.  Algorithmic
information and computable structures are developed systematically in
\cite{LiVitanyi2008,AshKnight2000}.  Complexity and descriptive complexity
are treated in \cite{Papadimitriou1994,AroraBarak2009,Immerman1999}.

The two lists have deliberately different standings.  A modern textbook can
supply a precise contemporary formulation without becoming the historical
source of the theorem.

%% Coverage: D-P5-007, N-a-064, A-050, A-052, A-127, A-186.
